\section{Normal kernels on the sextic Galois closure}\label{sec:heptic}
In degree seven the normal kernel at a ramified cubic has two real directions on its Galois closure. Projectivizing this plane does not separate the equations at the different nodes. We identify the one remaining real equation after solving the complex matching equations.
\subsection{The cubic and its Galois closure}\label{p14-sec:kernels}

Fix a complex rank-two bundle $N$ of degree $2g-2$ over a closed Riemann
surface $C$ of genus $g\ge2$. A real Cauchy--Riemann operator $D$ on $N$
has the Dolbeault symbol and a real-linear term of order zero. All kernels,
cokernels and indices below are real. Let $f:Y\to C$ be a connected cubic
with monodromy $S_3$ and $b>0$ distinct simple branch values. Write
$Z\to C$ for the smooth Galois closure and $V$ for the real plane
$\{(a_1,a_2,a_3):\sum a_i=0\}$ with its permutation action. Locally
the involution at a ramification point acts by $z\mapsto-z$, with
$w=z^2$ on the base. The fixed-target Hurwitz space of $f$ has complex
dimension $b$.

Let $D=D_0$ be invertible on $C$. At an \emph{ordinary transverse
cubic crossing} of a path $D_t$, the kernel of $D_f=f^*D$ is a line
$\mathbb R\xi$, the graph of $\xi$ in $N$ is embedded, and differentiation
at $\xi$ in the branch positions and increasing time gives an isomorphism
\begin{equation}\label{p14-eq:lowertransverse}
 W_f:T_f\mathcal H_f\oplus\mathbb R\longrightarrow
          \operatorname{coker}D_f .
\end{equation}
Indeed $\operatorname{ind}D_f=-2b$. The trace of $\xi$ is zero by
removability on $C$ and base invertibility. Its three conjugates on $Z$
therefore span one copy of $V$. Over a regular disk they give three
solutions with sum zero. Evaluation of this plane will have real rank
two at regular points and complex span of dimension two on at least
one regular disk. We explain both the generic assertion and a construction
on any prescribed $f$.

\begin{lemma}\label{p14-lem:kernelplane}
Suppose two of the sheet values form a complex frame on a regular disk.
For a residual set of arbitrarily small smooth variations on this disk
preserving the three solutions, every connected smooth finite cover
$u:X\to C$ has
\[
                     \ker u^*D=V^{K_u}.
\]
Here $K_u\le S_3$ is the image in the lower monodromy of a sheet
stabilizer for $X$ in the joint monodromy of $X$ and $Z$. The
identification uses pullback to a common normal cover.
\end{lemma}

\begin{proof}
Use real variations annihilating the two frame vectors $s_1,s_2$.
On any smaller disk a dense set of these variations attains the
infinite-order Petri condition at a specified interior point. This
condition says that a tensor of formal solution and adjoint-solution
jets with zero pointwise product jet is zero. For unrestricted
variations use \cite[Corollary~5.10]{wendl} and
\cite[Theorems~1.6.17 and~2.5.3]{dwbn}. For a nearby unrestricted
operator $D'$ with such a jet, solve $D's'_j=0$ locally with $s'_j$
close to $s_j$: a Cauchy right inverse on a disk of radius $r$ has
norm $O(r)$ on bounded functions, and interior elliptic estimates
give closeness in any prescribed finite smooth norm. Conjugate by the
complex gauge taking $s_j$ to $s'_j$, then cut off the difference from
$D$ away from the point. This preserves the Petri condition there and
annihilates the original columns on the entire disk.

Consider an irreducible real finite monodromy system $W$ inequivalent
to the prescribed standard system, with its ineffective punctures
removed. A nonzero twisted solution provides columns $w_i$ which,
together with $s_1,s_2$, are real linearly independent as germs on a
common normal cover. This follows from inequivalence, irreducibility
and unique continuation. An annihilator of the section and operator
variations has adjoint columns $v_i$ with
\[
        \sum_i w_i\otimes v_i=s_1\otimes\alpha_1+s_2\otimes\alpha_2.
\]
This identity is the pointwise annihilator condition for maps killing
the two prescribed columns. Test it against $[D,\Phi]$, with $\Phi$ a
compactly supported complex-linear endomorphism of $N$. The integral
on the left vanishes. The sections $\Phi s_1,\Phi s_2$ are
independently arbitrary, so both $\alpha_j$ solve the adjoint equation
on the disk. The Petri condition and independence of the input jets
force all $v_i$ to vanish to infinite order. Adjoint unique
continuation gives surjectivity.

If the $a$ effective punctures move, the projective incidence has
projection index
\[
       2a-1-2\sum_x\dim(W/W^{I_x})\le -1 ,
\]
by the orbifold index formula \cite[Proposition~2.8.6]{dwbn};
$I_x$ denotes inertia at $x$. Exhaust each configuration chart by
compact blocks whose punctures avoid a common smaller test disk.
At a good Petri jet over a block, the unit-normalized solution set is
compact by elliptic estimates on the compactified normal covers.
Finitely many allowed variations give surjectivity near that set.
Finite-dimensional Sard gives arbitrarily small perturbations excluding
all such solutions over the block. Absence is open there. There are
countably many blocks and discrete monodromy types. The remaining
type $V$ is faithful, so each of its copies on a common normal cover
descends to $Z$. Each has a one-dimensional invariant line on $Y$;
openness of the given simple kernel on the fixed $Y$ permits just one
copy. Taking stabilizer invariants proves the formula, by bounded
descent and removability over the branch points.
\end{proof}

At an ordinary crossing on a generic path the required regular-fiber
evaluation is injective everywhere. Indeed the two independent
differences of the three heights at a marked regular value can be
varied independently. For completeness, the universal section
equation at a cubic with complex frame disk admits arbitrary trace-free
forcing there by variation of $D$. In the adjoint test the trace
vanishes by base invertibility, hence the annihilator vanishes off
the marked points by unique continuation. Any remaining adjoint
distribution has finite order and is supported at the marks. Prescribe
holomorphic value and derivative jets there and solve recursively
for mixed jets so that the image under $D_f$ vanishes to higher
order than that distribution. Taylor polynomials with cutoffs realize
these jets. Independence of the resulting two value differences
kills their multipliers, and elliptic regularity removes the supported
adjoint. The same argument with a holomorphic
linear jet at the ramifying point gives independent variations of the
height difference at the two points in a branch fiber and the
ramification derivative. One may perform these tests without first
assuming a complex frame disk: two real-independent sheet values
suffice for the forcing test. Such a disk of real independence exists,
since collinear nonzero solutions on an open disk have constant real
ratio by their equation, contradicting the standard monodromy.
Real rank one for a pair of real four-vectors has codimension three,
and rank zero has codimension eight. The regular mark contributes
only two parameters along a transverse crossing. The complex
dependence locus at a fixed regular mark has codimension two.
These tests give the claimed ranks, including complex independence
on some disk. On continued transverse crossing charts the preceding
residual perturbations are available by multiplication by a time
cutoff.

If a nonzero section on a positive component is described by
$a\in V^{K_u}$, the component factors through
$Z/\operatorname{Stab}(a)$. Each nonzero stabilizer orbit has size
three or six. At a branch value the heights have values $u,u,-2u$
with $u\ne0$ by embeddedness. The kernel of evaluation is the
negative eigenline of inertia, generated, in a suitable ordering, by
$(1,-1,0)$. A nonzero vector in this line has trivial stabilizer.
Thus a positive component with a nonzero section vanishing somewhere
factors through $Z$.

In total degree seven some positive component must have zero height.
Connectedness forces a nonzero component to meet height zero since
contracted components carry constant sections. The degree count
leaves exactly one copy of $Z$ and one copy of $C$. Up to the deck
group their nonzero height on $Z$ is a scalar multiple of a fixed
root $s$, negated by a transposition $\tau$. Over each branch value
$x_i$ there is a unique ramification point $r_i$ of $Z$ fixed by
$\tau$. Its height is zero. The other two heights over $x_i$ are
nonzero and opposite. At a regular value the six heights are
distinct by real injectivity. At all three ramification points the
vertical derivatives are nonzero, being proportional to the
difference of two distinct coordinates of the root times the
ramification derivative of the lower embedded graph. The sextic
graph is therefore embedded. Its intersections with the zero section
are precisely the $b$ points just described, with vertical and
horizontal tangent lines respectively. The direct nodes or connected
contracted subcurves joining the two components determine a unique
nonempty subset $I$ of these points. No point is used twice since
the positive normalizations are embedded. Riemann--Hurwitz gives
$g(Z)=1+6(g-1)+3b/2$. Joining at $j$ points before attaching any
contracted components gives genus $1+7(g-1)+3b/2+j$.

\subsubsection{Construction on a fixed covering curve}
Here and in Section~\ref{p14-sec:joins}, prescribing the crossing permits
variation of the operator and its solution, with $C,N,f$ fixed.

\begin{proposition}\label{p14-prop:realize}
Every connected simply branched cubic $f$ as above admits an ordinary
transverse crossing with the stated evaluation ranks and central kernels.
Its branch weights can be chosen in any prescribed chamber of the
nonempty subset-sum hyperplanes.
\end{proposition}

\begin{proof}
Over the regular locus map the flat $V$-bundle real-linearly to
$N_{\mathbb R}$. At a branch send an even and an odd basis vector to
$e_0$ and $zv_0$, where $e_0,v_0$ form a complex basis.
This is equivariant under $z\mapsto-z$. Fix also a complex frame on
a regular disk. Extend relatively over the complement by a smooth
bundle map; transversality avoids the rank-one and rank-zero strata
of codimensions three and eight over the surface. The maps of the
three sheet vectors define an embedded cubic graph. Start with a
normal operator equal to $\bar\partial$ in the branch frames. Off
the branches prescribe a real zeroth-order correction on the image
of the flat plane by sending each parallel column to minus its
image under this operator. Extend smoothly on a complement, using
the zero correction near the branch centers. This makes the
prescribed heights solutions.

We verify regularity on the fixed $f$. On the frame disk variations
killing the two heights can first make $D$ invertible. A nonzero
base kernel vector cannot belong to their real span on an open
disk: its two real coefficients would have zero $\bar\partial$ and
be constant, contrary to the absence of an invariant in $V$.
Testing such a vector against a nonzero adjoint therefore gives
a nonzero kernel-to-cokernel matrix entry, by variation of $D$
and unique continuation. In a finite reduction the rank increases
under a small multiple of this variation. Iterate until the
index-zero base problem is invertible.

Augment $D_f$ by $b$ complex pole columns at the ramifying points.
Each column is the image off the pole of $\chi a/z$, where $a$ is
in the complex derivative line of the height there and $\chi$ is an
even cutoff. The pulled-back antilinear coefficient contains a
factor $\bar z$, so the forcing lies in $L^p$, $p>2$. The derivative
of the inverse sheet $(w-\alpha)^{1/2}$ is $-\alpha/(2z)$, giving
exactly the branch-motion column modulo the regular domain. The
augmented index is zero and the known kernel line persists. Suppose
an additional augmented solution exists. Its trace extends by
cancellation of the two odd poles at each ramifying pair and is zero
by invertibility. The adjoint of the augmented problem also has
trace zero by testing base pullbacks. Two of its three components
are real-independent on an open dense subset: proportionality on
an open set would give a constant real relation by the adjoint
equation and hence on the normal cover, whereas the nonzero
trace-free conjugates span the irreducible standard representation.
Writing the three conjugates of the extra solution as $h_i$, annihilation
of the constrained variations would put both $h_1-h_3,h_2-h_3$ in
the real span of $s_1,s_2$ there. Their coefficients would again be
constant. By trace zero and continuation all $h_i$ would lie in
the original standard copy and have no poles, so its sheet-invariant
line is the only possibility. Thus a rank-increasing variation exists
until the augmented kernel is just the prescribed line. One unrestricted
time variation on the disk then pairs nontrivially with the remaining
adjoint: the available forcing on the three sheet heights is any
triple of sum zero. This proves \eqref{p14-eq:lowertransverse}.
Lemma~\ref{p14-lem:kernelplane} applies in its open regular neighborhood.

Section~\ref{p14-sec:joins} below proves the assertion concerning strict
branch weights by changing the coefficients at the branch points and
restoring the equation on this same $Y$. This step also supplies the
additional regularity needed for the signed local degrees.
\end{proof}
